% Part: many-valued-logic
% Chapter: three-valued-logics
% Section: lukasiewicz

\documentclass[../../../include/open-logic-section]{subfiles}

\begin{document}

\olfileid[hi]{mvl}{thr}{luk}

\olsection{लुकासिएविच तर्कशास्त्र}

बहुमूल्य तर्कशास्त्र के आरंभिक प्रकाशित और सुविकसित प्रस्तावों में से एक
पोलिश दार्शनिक यान लुकासिएविच ने 1921 में प्रस्तुत किया था। लुकासिएविच को
अरस्तू की समुद्री युद्ध की समस्या ने प्रेरित किया था: ऐसा लगता है कि
\emph{आज} यह वाक्य, ``कल समुद्री युद्ध होगा'', न तो सत्य है और न असत्य:
उसका सत्य-मूल्य अभी निर्धारित नहीं हुआ है। लुकासिएविच ने ऐसे
``भविष्यगत आपातिक'' वाक्यों के लिए एक तीसरा सत्य-मूल्य जोड़ने का प्रस्ताव
रखा।
\begin{quote}
मैं बिना किसी विरोधाभास के यह मान सकता हूँ कि अगले वर्ष के किसी निश्चित
क्षण में, उदाहरणार्थ 21 दिसंबर को दोपहर के समय, वारसा में मेरी उपस्थिति
वर्तमान समय में न तो सकारात्मक रूप से और न नकारात्मक रूप से निर्धारित है।
अतः यह संभव है, किंतु आवश्यक नहीं, कि मैं दिए हुए समय पर वारसा में उपस्थित
रहूँगा। इस धारणा पर यह प्रतिज्ञप्ति, ``मैं अगले वर्ष 21 दिसंबर को दोपहर में
वारसा में रहूँगा'', वर्तमान समय में न तो सत्य हो सकती है और न असत्य। क्योंकि
यदि वह अभी सत्य होती, तो वारसा में मेरी भावी उपस्थिति आवश्यक होनी चाहिए थी,
जो धारणा के विरुद्ध है। दूसरी ओर, यदि वह अभी असत्य होती, तो वारसा में मेरी
भावी उपस्थिति असंभव होनी चाहिए थी, जो भी धारणा के विरुद्ध है। अतः विचाराधीन
प्रतिज्ञप्ति इस समय न सत्य है न असत्य और उसका कोई तीसरा मूल्य होना चाहिए, जो
``0'' अर्थात् असत्यता और ``1'' अर्थात् सत्यता से भिन्न हो। इस मूल्य को हम
``$\frac{1}{2}$'' से निर्दिष्ट कर सकते हैं। यह ``संभव'' को निरूपित करता है
और तीसरे मूल्य के रूप में ``सत्य'' तथा ``असत्य'' के साथ जुड़ता है।
\end{quote}
लुकासिएविच के तीसरे सत्य-मूल्य के लिए हम $\Undef$ का प्रयोग करेंगे।
\footnote{यहाँ लुकासिएविच ``संभव'' शब्द का प्रयोग आज के प्रचलित अर्थ से
भिन्न अर्थ में करते हैं, अर्थात् संभव किंतु आवश्यक नहीं।}

इस व्याख्या पर संयोजकों $\lnot$, $\land$ और $\lor$ के सत्य-फलन आसानी से
निर्धारित किए जा सकते हैं: किसी भविष्यगत आपातिक वाक्य का निषेध भी भविष्यगत
आपातिक वाक्य होता है, अतः $\tf{\lnot}(\Undef) = \Undef$। यदि किसी संयोजन
का एक संयोज्य अनिर्धारित और दूसरा सत्य हो, तो संयोजन भी अनिर्धारित होता
है---अंततः, भविष्यगत आपातिक संयोज्य का परिणाम कैसा निकलता है, उसके अनुसार
संयोजन सत्य भी निकल सकता है और असत्य भी। अतः
\[
    \tf{\land}(\True, \Undef) =
\tf{\land}(\Undef, \True) =
\Undef.
\]
परंतु यदि दूसरा संयोज्य असत्य हो, तो संयोजन सत्य नहीं निकल सकता, अतः
\[\tf{\land}(\False, \Undef) =
\tf{\land}(\False, \Undef) = \False.\]
अन्य मानों के लिए (जब तर्क निर्धारित सत्य-मूल्य, $\True$ या $\False$, हों)
परिणाम शास्त्रीय तर्कशास्त्र जैसे होते हैं।

सशर्त के मामले में स्थिति कुछ अधिक जटिल है। मान लें कि $q$ एक भविष्यगत
आपातिक कथन है। यदि $p$ असत्य है, तो $q$ का परिणाम चाहे जो हो, $p \lif q$
सत्य होगा, अतः हमें $\tf{\lif}(\False, \Undef) = \True$ रखना चाहिए। और यदि
$p$ सत्य है, तो $q$ का परिणाम चाहे जो हो, $q \lif p$ सत्य होगा, अतः
$\tf{\lif}(\Undef, \True) = \True$। यदि $p$ सत्य है, तो $p \lif q$ सत्य
या असत्य निकल सकता है, अतः $\tf{\lif}(\True, \Undef) = \Undef$। इसी तरह,
यदि $p$ असत्य है, तो $q \lif p$ सत्य या असत्य निकल सकता है, अतः
$\tf{\lif}(\Undef, \False) = \Undef$। अब वह स्थिति बचती है जिसमें $p$ और
$q$ दोनों भविष्यगत आपातिक हैं। प्रेरक विचार के आधार पर हमें इस स्थिति में
वास्तव में $\Undef$ देना चाहिए। किंतु इससे $!A \lif !A$ सर्वसत्यता
\emph{नहीं} रहेगा। लुकासिएविच को $!A \lor \lnot !A$ और
$\lnot(!A \land \lnot !A)$ छोड़ने में कठिनाई नहीं हुई, पर $!A \lif !A$
छोड़ने पर वे हिचक गए। इसलिए उन्होंने निर्धारित किया कि
$\tf{\lif}(\Undef, \Undef) = \True$।

\begin{defn}\ollabel{def:lukasiewicz}
त्रिमूल्य लुकासिएविच तर्कशास्त्र निम्नलिखित आव्यूह से परिभाषित होता है:
\begin{enumerate}
  \item $\lnot$, $\land$, $\lor$, $\lif$ वाली मानक प्रतिज्ञप्तिक भाषा
  $\Lang L_0$।
  \item सत्य-मूल्यों का समुच्चय $V = \{\True, \Undef, \False\}$।
  \item $\True$ एकमात्र निर्दिष्ट मूल्य है, अर्थात् $V^+ = \{\True\}$।
  \item सत्य-फलन निम्नलिखित सारणियों से दिए जाते हैं:
  \begin{center}
    \begin{tabular}{c|c} 
      $\tf{\lnot}$ & \\ 
      \hline  
      $\True$ & $\False$ \\ 
      $\Undef$ & $\Undef$ \\
      $\False$ & $\True$ 
    \end{tabular}
    \quad
    \begin{tabular}{c|ccc} 
      $\tf{\land}[\LogLuk[3]]$ & $\True$ & $\Undef$ & $\False$ \\ 
      \hline 
      $\True$ & $\True$ & $\Undef$ & $\False$ \\ 
      $\Undef$ & $\Undef$ & $\Undef$ & $\False$\\ 
      $\False$ & $\False$ & $\False$ & $\False$ 
    \end{tabular}
    \\[2ex]
    \begin{tabular}{c|ccc} 
      $\tf{\lor}[\LogLuk[3]]$ & $\True$ & $\Undef$ & $\False$ \\ 
      \hline 
      $\True$ & $\True$ & $\True$ & $\True$ \\ 
      $\Undef$ & $\True$ & $\Undef$ & $\Undef$ \\
      $\False$ & $\True$ & $\Undef$ & $\False$ 
    \end{tabular}
    \quad
    \begin{tabular}{c|ccc} 
      $\tf{\lif}[\LogLuk[3]]$ & $\True$ & $\Undef$ & $\False$ \\ 
      \hline 
      $\True$ & $\True$ & $\Undef$ & $\False$ \\ 
      $\Undef$ & $\True$ & $\True$ & $\Undef$  \\ 
      $\False$ & $\True$ & $\True$ & $\True$ 
    \end{tabular}
  \end{center} 
\end{enumerate}
\end{defn}

आसानी से देखा जा सकता है कि केवल $\lnot$, $\land$ और $\lor$ वाली कोई भी
!!{formula}~$!A$, यदि उसके सभी !!{propositional variable}ों को $\Undef$ दिया गया
हो, तो सत्य-मूल्य $\Undef$ लेगी। इसलिए, उदाहरण के लिए, शास्त्रीय
सर्वसत्यताएँ $p \lor \lnot p$ और $\lnot(p \land \lnot p)$
$\LogLuk[3]$ में सर्वसत्यताएँ नहीं हैं, क्योंकि $\pAssign v(p) = \Undef$
होने पर $\pValue{v}(!A) = \Undef$ होता है।

जिन !!{valuation} में $\pAssign v(p) = \True$ या $\False$ हो, उनमें
$\pValue v(!A)$ अपने शास्त्रीय सत्य-मूल्य के समान होगा।

\begin{prop}
  यदि $!A$ में आने वाले प्रत्येक $p$ के लिए $\pAssign v(p) \in
  \{\True, \False\}$, तो $\pValue v(!A)[\LogLuk[3]] =
  \pValue v(!A)[\LogCL]$।
\end{prop}

\begin{prob}\label{mvl:thr:luk:prob:luk-iff} मान लें कि हम $\LogLuk[3]$
  में $\pValue v(!A \liff !B) = \pValue v((!A \lif !B) \land
  (!B \lif !A))$ परिभाषित करते हैं। तब $\liff$ की सत्य-सारणी क्या होगी?
\end{prob}

बहुत-सी शास्त्रीय सर्वसत्यताएँ $\LogLuk[3]$ में भी सर्वसत्यताएँ
\emph{हैं}, जैसे $\lnot p \lif (p \lif q)$। शास्त्रीय तर्कशास्त्र की ही
तरह, हम सत्य-सारणी से इसकी जाँच कर सकते हैं:
\begin{center}
\begin{tabular}{cc|cccccc}
  $p$ & $q$ & $\lnot$ & $p$ & $\lif$ & $\smash{(}p$ & $\lif$ & $q\smash{)}$ \\ \hline
  \True & \True & \False & \True & \True & \True & \True & \True \\
  \True & \Undef & \False & \True & \True & \True & \Undef & \Undef \\
  \True & \False & \False & \True & \True & \True & \False & \False \\
  \Undef & \True & \Undef & \Undef & \True & \Undef & \True & \True \\
  \Undef & \Undef & \Undef & \Undef & \True & \Undef & \True & \Undef \\
  \Undef & \False & \Undef & \Undef & \True & \Undef & \Undef & \False \\
  \False & \True & \True & \False & \True & \False & \True & \True \\
  \False & \Undef & \True & \False & \True & \False & \True & \Undef \\
  \False & \False & \True & \False & \True & \False & \True & \False \\  
\end{tabular}
\end{center}

\begin{prob}
  दिखाइए कि निम्नलिखित $\LogLuk[3]$ में सर्वसत्यताएँ हैं:
  \begin{enumerate}
    \item $p \lif (q \lif p)$
    \item\label{mvl:thr:luk:prob:luk-taut-2} $\lnot(p \land q) \liff (\lnot p \lor \lnot q)$
    \item\label{mvl:thr:luk:prob:luk-taut-3} $\lnot(p \lor q) \liff (\lnot p \land \lnot q)$
  \end{enumerate}
  (\olref[mvl][thr][luk]{prob:luk-taut-2} और
  \olref[mvl][thr][luk]{prob:luk-taut-3} में $!A \liff !B$ को
  $(!A \lif !B) \land (!B \lif !A)$ का संक्षिप्त रूप मानें, या
  \olref[mvl][thr][luk]{prob:luk-iff} के अपने हल का उपयोग करें।)
\end{prob}

\begin{prob}
  दिखाइए कि निम्नलिखित शास्त्रीय सर्वसत्यताएँ $\LogLuk[3]$ में
  सर्वसत्यताएँ नहीं हैं:
  \begin{enumerate}
    \item $(\lnot p \land p) \lif q)$
    \item $((p \lif q) \lif p) \lif p$
    \item $(p \lif (p \lif q)) \lif (p \lif q)$
  \end{enumerate}
\end{prob}

इसलिए शायद यह सोचा जा सकता है कि यद्यपि सभी शास्त्रीय सर्वसत्यताएँ
$\LogLuk[3]$ में सर्वसत्यताएँ नहीं हैं, फिर भी प्रत्येक !!{valuation} पर
उन्हें कम-से-कम $\True$ या $\Undef$ मूल्य लेना चाहिए। किंतु ऐसा नहीं है।
इसका एक प्रतिदृष्टांत है
\[
  \lnot(p \lif \lnot p) \lor \lnot(\lnot p \lif p),
\]
जो $p$ का मूल्य $\Undef$ होने पर $\False$ है।

\begin{prob}
  निम्नलिखित में से कौन-से संबंध लुकासिएविच तर्कशास्त्र में मान्य हैं?
  प्रत्येक के लिए सत्य-सारणी दीजिए।
  \begin{enumerate}
    \item $p, p \lif q \Entails q$
    \item $\lnot\lnot p \Entails p$
    \item $p \land q \Entails p$
    \item $p \Entails p \land p$
    \item $p \Entails p \lor q$
  \end{enumerate}
\end{prob}

लुकासिएविच ने अपनी त्रिमूल्य प्रणाली के आधार पर संभावना का तर्कशास्त्र
बनाने की आशा की थी। इसके लिए उन्होंने एक-स्थानिक संयोजक $\Diamond !A$
(जिसका अर्थ है ``$!A$ संभव है'') और संगत $\Box !A$ (जिसका अर्थ है
``$!A$ आवश्यक है'') प्रस्तुत किए:
\begin{center}
  \begin{tabular}{c|c} 
    $\tf{\Diamond}$ & \\ 
    \hline  
    $\True$ & $\True$ \\ 
    $\Undef$ & $\True$ \\
    $\False$ & $\False$ 
  \end{tabular}
  \quad
  \begin{tabular}{c|c} 
    $\tf{\Box}$ & \\ 
    \hline  
    $\True$ & $\True$ \\ 
    $\Undef$ & $\False$ \\
    $\False$ & $\False$ 
  \end{tabular}
\end{center}
दूसरे शब्दों में, $p$ तभी संभव है जब वह पहले से असत्य के रूप में निर्धारित
न हो; और $p$ तभी आवश्यक है जब वह पहले से सत्य के रूप में निर्धारित हो।

\begin{prob}
  दिखाइए कि $\Box p \liff \lnot\Diamond \lnot p$ और $\Diamond p \liff
  \lnot \Box \lnot p$, $\Box$ और $\Diamond$ की सत्य-सारणियों से विस्तारित
  $\LogLuk[3]$ में सर्वसत्यताएँ हैं।
\end{prob}

किंतु इस प्रस्तावित मोडल तर्कशास्त्र की कमियाँ शीघ्र ही स्पष्ट हो गईं:
परिस्थितियाँ चाहे जैसी निकलें, $p \land \lnot p$ कभी सत्य नहीं निकल सकता।
इसलिए, भले ही अभी उसका निश्चय न हुआ हो (और इस कारण वह अनिर्धारित हो), उसे
असंभव माना जाना चाहिए; अर्थात् $\lnot \Diamond(p \land \lnot p)$ एक
सर्वसत्यता होनी चाहिए। किंतु यदि $\pAssign v(p) = \Undef$, तो
$\pValue v(\lnot \Diamond(p \land \lnot p)) = \Undef$। लुकासिएविच का यह
कहना सही था कि संभावना और आवश्यकता जैसे मोडल भेदों को समायोजित करने के लिए
दो सत्य-मूल्य पर्याप्त नहीं होंगे, पर तीसरा सत्य-मूल्य जोड़ना भी पर्याप्त
नहीं है।

\end{document}
